\section{Инфраструктура хранения и подготовки данных}

Настоящий раздел описывает технический контур, обеспечивающий хранение, нормализацию и аналитическую доступность данных iPinYou: от обоснования выбора СУБД и её развёртывания до определения схемы таблиц, формирования итоговой денормализованной таблицы и подключения к BI-инструменту Yandex DataLens.

\subsection{Развёртывание ClickHouse и создание таблиц}
\label{subsec:clickhouse_deploy}

\subsubsection*{Объём исходных данных и требования к хранилищу}

В работе использованы два сезона соревнования iPinYou (сезон 2 и сезон 3), имеющих идентичный формат. Совокупный объём загруженных файлов после предобработки приведён в Таблице~\ref{tab:file_sizes}.

\begin{table}[h!]
\small
  \begin{center}
    \caption{Объём обогащённых TSV-файлов по типам логов и сезонам.}
    \label{tab:file_sizes}
    \renewcommand{\arraystretch}{1.15}
    \begin{tabular}{l|r|r|r}
      \textbf{Тип лога} & \textbf{Сезон 2 (строк)} & \textbf{Сезон 3 (строк)} & \textbf{Итого (строк)} \\
      \hline
      Bidding log  & 53\,289\,330 & 11\,457\,419 & 64\,746\,749 \\
      Impression log & 12\,237\,087 &  3\,158\,171 & 15\,395\,258 \\
      Click log    &      9\,978  &      2\,716  &     12\,694  \\
      Conversion log &        494  &        539  &      1\,033  \\
    \end{tabular}
  \end{center}
\end{table}

Совокупный объём журнала ставок превышает 75 миллионов записей, а суммарный размер обогащённых TSV-файлов составляет порядка 30 ГБ. Данный масштаб данных определяет требования к хранилищу: необходимы эффективное сжатие, быстрое выполнение агрегирующих запросов по большим выборкам и поддержка вычисляемых столбцов.

\subsubsection*{Выбор СУБД: сравнение альтернатив}

При выборе инструмента хранения и обработки данных рассматривались три варианта.

\textbf{Загрузка файлов напрямую в Yandex DataLens.} DataLens поддерживает загрузку CSV/TXT-файлов в качестве источника данных. Однако данный подход неприменим в рамках настоящей работы по ряду причин: (1) любой запрос к файловому источнику выполняет полное сканирование всех строк без возможности использования индексов или ключей сортировки; (2) доступный набор функций и агрегаций существенно ограничен по сравнению с полноценной СУБД; (3) максимальный допустимый объём загружаемого файла составляет 400 МБ, тогда как совокупный объём данных настоящей работы превышает этот предел более чем в 70 раз.

\textbf{PostgreSQL.} Реляционная СУБД PostgreSQL~\cite{postgresql} является широко распространённым решением для хранения структурированных данных. Тем не менее для аналитических нагрузок, характерных для данной работы (агрегирование по десяткам миллионов строк, группировка по множеству измерений), PostgreSQL уступает колоночным СУБД: строковое хранение данных требует чтения всей строки даже при обращении к одному-двум столбцам, а отсутствие нативной векторизации запросов приводит к значительно более высокой латентности.

\textbf{ClickHouse.} Колоночная СУБД ClickHouse~\cite{clickhouse} оптимизирована для высокопроизводительной аналитической обработки (OLAP) и наилучшим образом соответствует требованиям настоящей работы. Ключевые преимущества ClickHouse применительно к задаче анализа RTB-логов:

\begin{itemize}
    \item \textbf{Колоночное хранение и сжатие.} Данные хранятся по столбцам, что позволяет читать только те столбцы, которые задействованы в запросе, и применять эффективные алгоритмы сжатия к однородным данным каждого столбца.
    \item \textbf{Векторизованное выполнение запросов.} Операции над данными выполняются блоками (batch processing), что позволяет эффективно использовать SIMD-инструкции процессора и существенно ускоряет агрегирование по большим таблицам.
    \item \textbf{Материализованные столбцы.} ClickHouse поддерживает объявление вычисляемых столбцов (\texttt{MATERIALIZED}), значения которых сохраняются на диске при вставке и не пересчитываются при каждом запросе (подробнее в подразделе~\ref{subsec:materialized_cols}).
    \item \textbf{Ключ сортировки и первичный индекс.} Движок \texttt{MergeTree} хранит данные отсортированными по заданному ключу и поддерживает разреженный первичный индекс, что обеспечивает эффективную фильтрацию по диапазонам времени и идентификаторам.
    \item \textbf{Нативный коннектор DataLens.} DataLens поддерживает прямое подключение к ClickHouse без промежуточных слоёв, что упрощает построение интерактивных дашбордов поверх аналитических запросов.
\end{itemize}

\subsubsection*{Установка и развёртывание}

ClickHouse устанавливался локально с использованием официальной документации~\cite{clickhouse-install}. 

После установки сервер запускается командой \texttt{clickhouse server}, а клиент -- командой \texttt{clickhouse client}. Загрузка данных осуществлялась путём прямой передачи предварительно подготовленных TSV-файлов в стандартный поток ввода.

\subsubsection*{Схема таблиц и выбор движка}

Для каждого типа лога создана отдельная таблица: \texttt{bid\_log}, \texttt{imp\_log}, \texttt{clk\_log}, \texttt{conv\_log}. Все таблицы используют движок \texttt{MergeTree} с сортировкой по вычисляемому столбцу \texttt{ts}. Движок \texttt{MergeTree} выбран как базовый вариант для однократно загружаемых аналитических данных без необходимости дедупликации: он обеспечивает эффективное хранение и быстрое сканирование диапазонов по ключу сортировки.

Ключевые решения по типам данных:

\begin{itemize}
    \item \textbf{Timestamp} хранится в исходном виде как \texttt{String} в поле \texttt{timestamp\_raw} (формат \texttt{yyyyMMddHHmmssSSS}); разбор в тип \texttt{DateTime64(3)} выполняется материализованным столбцом \texttt{ts} (подробнее в подразделе~\ref{subsec:materialized_cols}).
    \item \textbf{Ad Slot Width / Height} объявлены как \texttt{Nullable(UInt32)}: поля могут содержать пустые значения, а в исходных данных встречаются записи с отсутствующими размерами слота.
    \item \textbf{Ad Exchange} объявлен как \texttt{UInt8 DEFAULT 0}: значение 0 соответствует введённой при предобработке категории «Unknown» для записей с отсутствующим идентификатором биржи.
    \item \textbf{Advertiser ID} в исходных таблицах объявлен как \texttt{Nullable(UInt32)}; в итоговой аналитической таблице (см. подраздел~\ref{subsec:ads_data}) переведён в обязательное поле с \texttt{DEFAULT 0} для обеспечения корректной сортировки и группировки.
    \item \textbf{Platform / Browser} -- два дополнительных столбца, добавленных скриптом обогащения на этапе предобработки; хранятся как \texttt{String}.
\end{itemize}

Таблица \texttt{imp\_log} дополнительно содержит поля \texttt{paying\_price} (\texttt{Nullable(UInt32)}) и \texttt{landing\_page\_url} (\texttt{Nullable(String)}), отсутствующие в \texttt{bid\_log}. Таблицы \texttt{clk\_log} и \texttt{conv\_log} имеют схему, идентичную \texttt{imp\_log}.

Полные DDL-определения всех четырёх таблиц приведены в Приложении~\ref{app:ddl}.

\subsection{Материализованные столбцы}
\label{subsec:materialized_cols}

В ClickHouse механизм \texttt{MATERIALIZED} позволяет объявить столбец, значение которого вычисляется один раз при вставке строки и сохраняется на диске наравне с остальными данными. Это принципиально отличается от вычисления выражения в каждом запросе: при аналитических сканированиях по десяткам миллионов строк однократная материализация устраняет повторные вычисления и позволяет использовать столбец в ключе сортировки, что ускоряет диапазонные запросы по времени.

\subsubsection*{Материализованный столбец \texttt{ts}}

Исходный формат метки времени в данных iPinYou -- строка вида \texttt{yyyyMMddHHmmssSSS} (17 символов, без разделителей). Поскольку ClickHouse не разбирает данный формат нативно, столбец \texttt{ts} объявлен как \texttt{MATERIALIZED} с выражением, которое при помощи функции \texttt{substring} извлекает компоненты даты и времени, конкатенирует их в стандартный формат \texttt{YYYY-MM-DD HH:MM:SS.mmm} и передаёт в \texttt{parseDateTime64BestEffort} с точностью до миллисекунды (\texttt{DateTime64(3)}). Результирующий столбец \texttt{ts} используется как ключ сортировки (\texttt{ORDER BY ts}) во всех четырёх таблицах, что обеспечивает эффективную фильтрацию по временным диапазонам.

\subsubsection*{Обогащение признаками \texttt{platform} и \texttt{browser}}

Разбор поля \texttt{user\_agent} в признаки платформы и браузера выполнен на этапе предобработки Python-скриптом, а не на уровне СУБД. Скрипт принимает нормализованный TSV-файл и дописывает в конец каждой строки два дополнительных поля: \texttt{platform} и \texttt{browser}. Таким образом, в ClickHouse эти столбцы поступают уже как обычные строковые поля, а не как \texttt{MATERIALIZED}-выражения.

Логика определения платформы: \texttt{Android}, \texttt{iPhone} (включает iPad и iPod), \texttt{Windows}, \texttt{Mac}, \texttt{Linux}, \texttt{Other}. Логика определения браузера: \texttt{IE} (MSIE/Trident), \texttt{Chrome}, \texttt{Firefox}, \texttt{UCBrowser}, \texttt{Baidu}, \texttt{Safari}, \texttt{Other}. При отсутствии или некорректном значении \texttt{user\_agent} оба поля принимают значение \texttt{Unknown}.
Полный исходный код скрипта приведён в Приложении~\ref{app:enricher}.

Подобный подход -- предварительное обогащение на уровне файла, а не вычисление в каждом SQL-запросе -- оправдан при однократной загрузке данных: стоимость разбора строки User-Agent значительно выше, чем чтение уже готового строкового значения из колонки.

\subsection{Формирование итоговой аналитической таблицы}
\label{subsec:ads_data}

\subsubsection*{Логика объединения таблиц}

Итоговая аналитическая таблица \texttt{ads\_data} формируется путём объединения четырёх исходных таблиц по ключу \texttt{bid\_id}. Центральной таблицей является \texttt{bid\_log}, содержащая все зафиксированные ставки DSP. К ней последовательно присоединяются:

\begin{itemize}
    \item \texttt{imp\_log} -- для получения факта победы в аукционе и фактической цены показа (\texttt{paying\_price});
    \item \texttt{clk\_log} -- для получения факта клика по объявлению и URL целевой страницы (\texttt{landing\_page\_url});
    \item \texttt{conv\_log} -- для получения факта конверсии.
\end{itemize}

Все соединения выполняются как \texttt{ANY LEFT JOIN}: оператор \texttt{ANY} в ClickHouse при наличии нескольких совпадений по ключу берёт первую найденную строку, что корректно для данной задачи, поскольку каждому \texttt{bid\_id} соответствует не более одного события каждого типа. Семантика \texttt{LEFT JOIN} гарантирует сохранение всех строк из \texttt{bid\_log}, в том числе проигравших ставок, для которых нет соответствующих записей в остальных таблицах.

По результатам объединения в таблицу добавляются три булевых флага:

\begin{itemize}
    \item \texttt{is\_win} -- ставка выиграла аукцион (присутствует запись в \texttt{imp\_log});
    \item \texttt{is\_clicked} -- показ сопровождался кликом пользователя (присутствует запись в \texttt{clk\_log});
    \item \texttt{is\_conversion} -- клик привёл к конверсии (присутствует запись в \texttt{conv\_log}).
\end{itemize}

Флаги вычисляются функцией \texttt{notEmpty(bid\_id)} из соответствующей присоединённой таблицы: если соединение не нашло совпадения, поле \texttt{bid\_id} будет пустой строкой, и \texttt{notEmpty} вернёт 0.

Полный SQL-запрос формирования таблиц \texttt{joined\_ads\_data} и \texttt{ads\_data} приведён в Приложении~\ref{app:merge_sql}.

\subsubsection*{Итоговая таблица \texttt{ads\_data}}

Поскольку поле \texttt{advertiser\_id} в промежуточной таблице объявлено как \texttt{Nullable(UInt32)}, а большинство аналитических запросов предполагают группировку и фильтрацию по рекламодателю, была выполнена дополнительная трансформация: пустые значения \texttt{advertiser\_id} заменены на 0 (категория «Unknown»), тип поля переведён в обязательный \texttt{UInt32}, и таблица пересоздана с ключом сортировки \texttt{ORDER BY advertiser\_id}. Результирующая таблица получила название \texttt{ads\_data} и является основным источником данных для всех последующих аналитических запросов и визуализаций.

\subsubsection*{Преимущества денормализованного подхода}

Денормализованная «широкая» таблица \texttt{ads\_data} предпочтительна для аналитических нагрузок в ClickHouse по ряду причин. Во-первых, ClickHouse оптимизирован для сканирования больших объёмов данных в одной таблице, тогда как выполнение JOIN в реальном времени на десятках миллионов строк существенно увеличивает латентность запросов. Во-вторых, хранение всех атрибутов события в одной строке позволяет DataLens и SQL-запросам обращаться к любому срезу (по времени, рекламодателю, платформе, домену) без дополнительных шагов. В-третьих, однократно вычисленные флаги \texttt{is\_win}, \texttt{is\_clicked}, \texttt{is\_conversion} устраняют необходимость повторного выполнения JOIN при каждом расчёте воронки или метрик эффективности.

\subsubsection*{Подключение ClickHouse к DataLens}
\label{subsec:datalens_dataset}

Для обеспечения доступа к локальному серверу ClickHouse из облачного сервиса Yandex DataLens применялся инструмент \texttt{ngrok}~\cite{ngrok}, осуществляющий проброс локального порта через защищённый туннель. Подключение DataLens к базе данных настроено с использованием TLS-сертификата и парольной аутентификации, что обеспечивает минимальные требования безопасности при работе с публично доступным туннелем.
